\chapter{LITERATURE REVIEW}

\section{General Overview}
The response of a building during an earthquake depends not only on the superstructure but also on the behaviour of the foundation and supporting ground. Wolf's formulation of dynamic soil--structure interaction established the coupled nature of these components and the limitations of treating the foundation as perfectly fixed \cite{wolf1985}. More recent studies continue to show that soil flexibility can alter vibration characteristics, deformation demand, and force distribution, particularly when the supporting ground is relatively soft \cite{poudel2024,kc2025}.

SSI is commonly discussed in terms of kinematic and inertial interaction. Kinematic interaction concerns the way a foundation modifies the motion that would exist in the free field, while inertial interaction arises from forces generated by the vibrating structure and transmitted through the foundation to the soil \cite{wolf1985,kc2025}. Both mechanisms are relevant to performance assessment because they can change the response quantities used to judge serviceability and damage, including displacement and inter-storey drift.

\section{Seismic Vulnerability of Buildings in Nepal}
Nepal's seismic vulnerability reflects both regional tectonic activity and characteristics of the building stock. Gautam et al. identified recurring structural and construction deficiencies in Nepalese buildings, while studies following the 2015 Gorkha earthquake documented the consequences of these weaknesses during strong shaking \cite{gautam2016,adhikari2022}. Paudel et al. likewise examined vulnerability in Nepalese reinforced-concrete residential buildings and highlighted the sensitivity of seismic performance to structural configuration \cite{paudel2024seismic}.

Local ground conditions add another source of uncertainty. Deep alluvial deposits, particularly in the Kathmandu Valley, can modify earthquake motion and contribute to differences between fixed-base predictions and the response of structures supported on deformable soil \cite{poudel2024,adhikari2022}. Poudel and Chaulagain's nonlinear investigation and the simplified study by Timilsina et al. both indicate that soil flexibility can materially affect displacement and drift demand in reinforced-concrete buildings \cite{poudel2024,timilsina2023}.

\section{Hospital Buildings as Critical Facilities}
The importance of hospital seismic performance in Nepal has been recognized for more than two decades. Guragain and Dixit's assessment of Nepalese hospitals showed that many facilities required significant improvement to achieve acceptable earthquake performance \cite{guragain2004}. More recent work by Subedi compared seismic demand in an RC hospital building using IS 1893:2016 and NBC 105:2020, illustrating how design provisions can influence calculated response \cite{subedi2022}.

These studies are especially relevant because the acceptable outcome for a hospital is not simply avoidance of collapse. The building must retain sufficient structural and functional capacity to support emergency services after an earthquake. This requirement makes deformation-sensitive measures, foundation response, and the distribution of structural demand important alongside conventional force-based checks \cite{guragain2004,subedi2022}.

\section{Soil--Structure Interaction Theory}
Wolf describes SSI as a dynamic interaction problem in which the motion of the foundation and the response of the superstructure influence each other through the supporting medium \cite{wolf1985}. On stiff ground, the fixed-base idealization may provide a reasonable approximation; as soil compliance increases, foundation translation and rotation become part of the total structural response. KC et al. summarize these effects within a broader review of ground-motion modification by soil, structures, and topography \cite{kc2025}.

The practical consequence is a change in effective stiffness, period, damping, and deformation pattern. Santisi d'Avila and Lopez-Caballero, for example, examined nonlinear interaction in a coupled soil--structure system, while Adhikari et al. demonstrated the importance of site condition in fragility analysis of RC buildings on alluvial deposits \cite{santisi2016,adhikari2022}. These findings support the use of flexible-base models when the study objective is to compare the response of the same building across different soil conditions.

\subsection*{Kinematic and inertial contributions to SSI}
Kinematic interaction arises when the presence and geometry of a foundation modify the motion transmitted from the free field to its base. Inertial interaction, by contrast, develops when earthquake-induced structural forces are transferred through the foundation to the supporting soil. The two mechanisms need not produce identical effects in the horizontal and vertical directions, and they may affect both the effective input motion and the foundation restraint \cite{wolf1985,kc2025}. The discrete support-spring comparison in this study focuses principally on changes in foundation restraint; it is not a replacement for an explicit calculation of free-field ground-motion modification.

\subsection*{Period lengthening and redistribution of structural demand}
The period of an equivalent single-degree-of-freedom system varies approximately with the square root of the ratio of participating mass to lateral stiffness. This conceptual relationship explains why adding foundation compliance commonly increases the vibration period. In a multi-degree-of-freedom building, however, the response depends on the distribution of masses, stiffnesses, modes and excitation. Response-spectrum force and displacement changes must therefore be evaluated independently rather than inferred from the sign of period change \cite{wolf1985,timilsina2023,poudel2024}.

\section{Foundation Spring and Flexible-Base Modelling}
SSI may be represented either through direct modelling of the soil domain or through a substructure approach in which foundation impedances are introduced into the structural model \cite{wolf1985,gajan2008}. The substructure approach is attractive for building analysis because translational and rotational stiffnesses can be defined at the foundation level while the superstructure remains within a conventional structural-analysis program.

The classical Winkler concept idealizes the supporting soil as a set of independent linear springs for which the local reaction is proportional to the local foundation displacement. In its simplest vertical form this is written as $p=k_s w$, where $p$ is contact pressure, $w$ is local displacement, and $k_s$ is the modulus of subgrade reaction. Contemporary SSI reviews describe the Winkler model as a simple spring representation of the soil domain, while also noting its limitations for reproducing the continuum response of soil \cite{krishnan2024}. The present study uses this spring-idealization concept at the structural-model level, but it does not obtain the six foundation stiffness components from a single subgrade-reaction modulus.

A common theoretical basis for selecting directional spring stiffnesses is the elastic-half-space foundation-impedance formulation assembled by Gazetas. Gazetas provides algebraic expressions and charts for surface and embedded rigid foundations, including vertical, horizontal, rocking, and torsional modes, as functions of soil properties and foundation geometry \cite{gazetas1991}. For this study, translational, rocking and torsional stiffness coefficients were assigned to the discrete foundation spring representation for each soil-support condition. Gajan et al. and Raychowdhury extend related spring concepts to nonlinear shallow-foundation mechanisms such as uplift, gapping, sliding, and cyclic settlement \cite{gajan2008,raychowdhury2008}; those nonlinear foundation mechanisms are outside the present linear spring idealization.

\subsection*{Winkler and Gazetas foundation idealizations}
In the Winkler representation, the soil below a foundation is idealized as a distribution of independent reaction springs. If the local vertical settlement is $w$ and the local contact pressure is $p$, the elastic relation is
\begin{equation}
p=k_s w,
\label{eq:winkler_relation}
\end{equation}
\RegisterThesisEquation{Winkler pressure--settlement relation}
where $k_s$ is the modulus of subgrade reaction, with units of force per unit volume when $p$ and $w$ are in consistent units. For a discretized foundation, a local spring supporting tributary area $A_i$ has the idealized vertical stiffness $k_{z,i}=k_s A_i$. This relationship expresses the simplest local compression response but does not by itself define the coupled horizontal, rocking or torsional impedance of a rigid footing \cite{krishnan2024,wolf1985}.

Gazetas (1991) presented dimension-dependent foundation impedance expressions for rigid foundations on or embedded in an elastic half-space \cite{gazetas1991}. The translational terms resist horizontal and vertical movement, the rocking terms resist rotations about horizontal axes, and the torsional term resists rotation about the vertical axis. For a general uncoupled spring representation, the resulting stiffness components are written as $K_x$, $K_y$, $K_z$, $K_{\theta x}$, $K_{\theta y}$ and $K_t$. Their numerical magnitudes depend on the relevant footing dimensions and soil shear modulus, Poisson's ratio and other modelling choices. A rigorous derivation must use the appropriate geometry and formulation rather than treating all components as multiples of the Winkler vertical spring.

\begin{figure}[htbp]
\centering\includegraphics[width=0.90\textwidth]{figures/ch2/winkler_gazetas_schematic.pdf}
\caption{Contrasting the local Winkler spring idealization with equivalent translational, rocking and torsional foundation restraint (conceptual schematic).}
\label{fig:winkler_literature}
\end{figure}

\subsection*{Modelling assumptions and applicability}
The use of independent linear springs is computationally convenient, but real soils can exhibit frequency dependence, material nonlinearity, uplift, redistribution of contact pressure and permanent settlement. Elastic impedances should not be interpreted as a complete representation of these processes. Studies on shallow-foundation response explicitly address mechanisms such as gapping, uplift and cyclic yielding that are outside an uncoupled elastic-spring representation \cite{gajan2008,raychowdhury2008}. The present comparisons consequently describe the influence of the spring properties assigned to the analysed models and are not independent determinations of site-specific nonlinear foundation behaviour.

\section{Influence of Soil Condition on Seismic Response}
The effect of SSI is not uniformly beneficial or adverse. Period lengthening can shift a structure to a different region of the response spectrum, which may reduce some force measures while increasing displacement or drift. Timilsina et al. reported noticeable changes in seismic response when soil flexibility was introduced using a simplified method, and Krishnan and Sivakumar similarly emphasized the influence of SSI on the stability and response of RC structures \cite{timilsina2023,krishnan2024}.

Poudel and Chaulagain investigated nonlinear SSI effects in irregular RC buildings and found that the interaction between soil and structure can substantially modify response compared with a fixed-base representation \cite{poudel2024}. Taken together, these studies indicate that the direction and magnitude of SSI effects should be established from the response quantities of the specific model rather than assumed in advance.

\section{Performance-Based Seismic Assessment}
Performance-based assessment evaluates seismic demand in relation to defined performance objectives rather than relying solely on elastic force checks. FEMA 356 provides a widely used framework for performance levels such as Immediate Occupancy, Life Safety, and Collapse Prevention \cite{fema356}. The approach places particular emphasis on deformation and component behaviour, which makes it compatible with an SSI study where foundation flexibility can alter drift and displacement.

Janous and El Ghoulbzouri examined the vulnerability of irregular RC buildings while considering SSI, and more recent studies have continued to incorporate soil effects into performance-oriented evaluation \cite{janous2024,oggu2025}. For hospitals, these methods are relevant because the expected post-earthquake condition is tied to continued service, not only to collapse resistance \cite{subedi2022}.

\section{Fragility Analysis and Damage States}
Fragility curves express the probability that a defined damage state will be exceeded at a specified intensity measure. Lognormal models are commonly used to relate intensity measures such as peak ground acceleration or spectral acceleration to damage-state exceedance probability \cite{giordano2021,adhikari2022}. The choice of engineering demand parameter and damage thresholds therefore has a direct influence on the resulting fragility relationship.

For Nepalese buildings, Giordano et al. developed empirical fragility models for school buildings, while Adhikari et al. analysed RC buildings with brick infills in a high-seismic region with alluvial deposits \cite{giordano2021,adhikari2022}. Kircil and Ethemoglu and Oggu and Gopikrishna examined the effect of SSI on damage probability and fragility in multistorey RC frames \cite{kircil2025,oggu2025}. These studies provide the theoretical context for understanding fragility curves and drift-based response classifications in seismic performance studies.

\subsection*{Deterministic response, damage states and probabilistic fragility}
Displacement, drift and member forces obtained from elastic response-spectrum analysis are engineering demand parameters, but they are not in themselves probabilities of damage-state exceedance. A fragility curve requires a specified intensity measure, damage state, response model and statistical description of demand or capacity. A widely used representation is the cumulative lognormal function
\begin{equation}
P(DS\geq ds_i\mid IM=x)=\Phi\!\left(\frac{\ln(x/\theta_i)}{\beta_i}\right),
\label{eq:fragility_lognormal}
\end{equation}
\RegisterThesisEquation{Lognormal damage-state exceedance probability}
where $\theta_i$ is the median intensity at exceedance of damage state $ds_i$, $\beta_i$ is the logarithmic dispersion, and $\Phi$ denotes the standard normal cumulative distribution. The formula is given as theoretical context; it is not asserted that these parameters have been independently identified for the submitted hospital-block curves \cite{giordano2021,adhikari2022}.

A smooth visual interpolation through tabulated fragility values is different from estimating a new physical or probabilistic model. The former preserves known ordinates; the latter changes the model's parameters and requires a documented fitting and validation procedure. The damage-state definitions associated with each fragility curve should be distinguished from the adopted inter-storey-drift screening bands.

\section{Research Gap}
The literature establishes that SSI can influence the seismic response of reinforced-concrete buildings, but Nepal-focused work on mid-size hospital buildings remains limited relative to studies of residential, school, or generic RC frames \cite{adhikari2022,paudel2024seismic,giordano2021,yadav2023}. Existing hospital studies emphasize vulnerability and code-based demand, whereas fewer studies integrate flexible-base modelling with response-spectrum demand, foundation response, member-force redistribution, and seismic fragility in hospital buildings \cite{guragain2004,subedi2022}.

The present thesis examines the six structural blocks of the G+4 Infectious and Communicable Disease Hospital case study in Lekhnath, Pokhara under fixed, hard-soil, medium-soil, and soft-soil support conditions. The primary comparison uses the block-specific modal periods and mass participation data, and compares displacement, drift and fragility across the foundation-support conditions. The arrangement permits a consistent within-block assessment of support flexibility.

\section{Summary of Literature Review}
The reviewed studies show that SSI changes the dynamic system rather than producing a single predictable change in seismic demand. Period, displacement, drift, shear, foundation response, and member demand can respond differently as soil flexibility increases. This is especially important for hospitals, where deformation and post-earthquake functionality are central to performance assessment.

The literature also supports drift-based fragility methods for connecting seismic demand with performance interpretation. For the present study, these findings support a model-by-model comparison of the four support conditions and provide the conceptual basis for interpreting variations in modal and seismic response.
